% Part: model-theory
% Chapter: models-of-arithmetic
% Section: non-standard-models

\documentclass[../../../include/open-logic-section]{subfiles}

\begin{document}

\olfileid{mod}{mar}{nst}
\section{النماذج غير القياسية}

\begin{explain}
معيار القياسية هو التشاكل: تكون !!a{structure} للغة
$\Lang{L_A}$ قياسية متى تشاكلت مع~$\Struct{N}$؛ وأما
!!a{structure} التي لا تتشاكل مع~$\Struct{N}$ فهي غير قياسية.
\end{explain}

\begin{defn}
تكون !!{structure}~$\Struct{M}$ للغة $\Lang{L_A}$ \emph{غير قياسية}
إذا لم تكن متشاكلة مع~$\Struct{N}$. وتسمى الـ!!{element}s
$\OLId{latin-ordinary}{x} \in \Domain{\OLId{latin-looped}{M}}$ التي تساوي $\Value{\num{\OLId{latin-ordinary}{n}}}{\OLId{latin-looped}{M}}$ لبعض
$\OLId{latin-ordinary}{n} \in \Nat$ \emph{أعدادًا قياسية} (في $\Struct{M}$)، وما عداها
\emph{أعدادًا غير قياسية}.
\end{defn}

\begin{explain}
تقتضي \olref[stm]{prop:standard-domain} أن كل !!{structure}
قياسية للغة~$\Lang{L_A}$ لا تضم إلا !!{element}s قياسية.
لكن العكس لا يلزم على العموم: فقد تكون الـ!!{structure}
غير قياسية ومجالها مستنفدًا بالأعداد القياسية، كما طُلب
في تمرين القسم السابق. وإنما الذي يلزم هو أن وجود
!!{element} غير قياسي يمنع قياسية الـ!!{structure}.
\end{explain}

\begin{prop}
إذا احتوت !!a{structure}~$\Struct{M}$ للغة $\Lang{L_A}$ عددًا غير
قياسي، كانت $\Struct{M}$ غير قياسية.
\end{prop}

\begin{proof}
لو كانت $\Struct{M}$ قياسية وفيها عدد غير قياسي~$\OLId{latin-ordinary}{x}$،
لوجد تشاكل $\OLId{latin-ordinary}{g}\colon \Nat \to \Domain{\OLId{latin-looped}{M}}$. وباستقراء يسير
على~$\OLId{latin-ordinary}{n}$ نحصل على
$\OLId{latin-ordinary}{g}(\Value{\num{\OLId{latin-ordinary}{n}}}{\OLId{latin-looped}{N}})=\Value{\num{\OLId{latin-ordinary}{n}}}{\OLId{latin-looped}{M}}$. فالدالة $\OLId{latin-ordinary}{g}$
ترسل الأعداد القياسية في~$\Struct{N}$ إلى القياسية في~$\Struct{M}$.
وحينئذ يبقى العدد غير القياسي في $\Struct{M}$ خارج صورتها،
فتخالف $\OLId{latin-ordinary}{g}$ شرط !!{surjective}، وهذا تناقض.
\end{proof}

\begin{prob}
تذكر أن $\Th{Q}$ تحتوي البديهيات
\begin{align*}
& \lforall[\OLId{latin-ordinary}{x}][\lforall[\OLId{latin-ordinary}{y}][(\eq[\OLId{latin-ordinary}{x}'][\OLId{latin-ordinary}{y}'] \lif \eq[\OLId{latin-ordinary}{x}][\OLId{latin-ordinary}{y}])]] \tag{$!Q_\OLMathNumeral{1}$}\\
& \lforall[\OLId{latin-ordinary}{x}][\eq/[\Obj 0][\OLId{latin-ordinary}{x}']] \tag{$!Q_\OLMathNumeral{2}$}\\
& \lforall[\OLId{latin-ordinary}{x}][(\eq[\OLId{latin-ordinary}{x}][\Obj 0] \lor \lexists[\OLId{latin-ordinary}{y}][\eq[\OLId{latin-ordinary}{x}][\OLId{latin-ordinary}{y}']])] \tag{$!Q_\OLMathNumeral{3}$}
\end{align*}
أعط !!{structure}s~$\Struct{M_1}$ و$\Struct{M_2}$ و$\Struct{M_3}$ بحيث
\begin{enumerate}
\item $\Sat{\OLId{latin-looped}{M}_\OLMathNumeral{1}}{!Q_\OLMathNumeral{1}}$ و$\Sat{\OLId{latin-looped}{M}_\OLMathNumeral{1}}{!Q_\OLMathNumeral{2}}$ و$\Sat/{\OLId{latin-looped}{M}_\OLMathNumeral{1}}{!Q_\OLMathNumeral{3}}$؛
\item $\Sat{\OLId{latin-looped}{M}_\OLMathNumeral{2}}{!Q_\OLMathNumeral{1}}$ و$\Sat/{\OLId{latin-looped}{M}_\OLMathNumeral{2}}{!Q_\OLMathNumeral{2}}$ و$\Sat{\OLId{latin-looped}{M}_\OLMathNumeral{2}}{!Q_\OLMathNumeral{3}}$؛ و
\item $\Sat/{\OLId{latin-looped}{M}_\OLMathNumeral{3}}{!Q_\OLMathNumeral{1}}$ و$\Sat{\OLId{latin-looped}{M}_\OLMathNumeral{3}}{!Q_\OLMathNumeral{2}}$ و$\Sat{\OLId{latin-looped}{M}_\OLMathNumeral{3}}{!Q_\OLMathNumeral{3}}$.
\end{enumerate}
ويكفي في كل مثال تعيين~$\Assign{\Obj{0}}{\OLId{latin-looped}{M}_\OLId{latin-ordinary}{i}}$ و
$\Assign{\prime}{\OLId{latin-looped}{M}_\OLId{latin-ordinary}{i}}$ لكل منها.
\end{prob}

\begin{explain}
أما إنشاء !!{structure}s غير قياسية للغة $\Lang{L_A}$
فيسير: نأخذ الـ!!{domain}~$\Int$، ونؤول الرموز غير المنطقية
بالمعاني المعتادة. ففيه أعداد سالبة لا تساوي قيمة
$\num{\OLId{latin-ordinary}{n}}$ لأي~$\OLId{latin-ordinary}{n}$، فالبنية غير قياسية. لكنها ليست
\emph{نموذجًا} للحساب بمعنى الاتفاق مع~$\Struct{N}$ في الجمل
الصادقة؛ فإن $\lforall[\OLId{latin-ordinary}{x}][\eq/[\OLId{latin-ordinary}{x}'][\Obj{0}]]$ تكذب فيها،
مثلًا. وأما وجود نماذج غير قياسية توافق الحساب نفسه،
فتثبته مبرهنة التراص بسهولة.
\end{explain}

\begin{prop}
لتكن $\Th{TA} = \Setabs{!A}{\Sat{\OLId{latin-looped}{N}}{!A}}$ نظرية~$\Struct{N}$.
لـ$\Th{TA}$ نموذج غير قياسي !!a{enumerable}.
\end{prop}

\begin{proof}
نزيد $\Lang{L_A}$~!!{constant} جديدًا~$\OLId{latin-ordinary}{c}$، ونأخذ
مجموعة الـ!!{sentence}s
\[
\OLId{greek-variable}{Gamma} = \Th{TA} \cup \{\eq/[\OLId{latin-ordinary}{c}][\num{\OLMathNumeral{0}}], \eq/[\OLId{latin-ordinary}{c}][\num{\OLMathNumeral{1}}],
\eq/[\OLId{latin-ordinary}{c}][\num{\OLMathNumeral{2}}], \dots\}.
\]
أي نموذج~$\Struct{M^c}$ لـ$\OLId{greek-variable}{Gamma}$ يشتمل على
!!a{element}~$\OLId{latin-ordinary}{x} =
\Assign{\OLId{latin-ordinary}{c}}{\OLId{latin-looped}{M}}$ غير قياسي، إذ $\OLId{latin-ordinary}{x} \neq \Value{\num{\OLId{latin-ordinary}{n}}}{\OLId{latin-looped}{M}}$
لكل $\OLId{latin-ordinary}{n} \in \Nat$. ولدينا أيضًا $\Sat{\OLId{latin-looped}{M}^\OLId{latin-ordinary}{c}}{\Th{TA}}$،
لأن $\Th{TA} \subseteq \OLId{greek-variable}{Gamma}$. فإذا رددنا $\Struct{M^c}$
إلى !!a{structure}~$\Struct{M}$ للغة $\Lang{L_A}$ بإسقاط~$\OLId{latin-ordinary}{c}$
من التأويل، بقي $\OLId{latin-ordinary}{x}$ غير القياسي في المجال، وبقي
$\Sat{\OLId{latin-looped}{M}}{\Th{TA}}$. وإنما يبقى هذا الصدق لأن $\OLId{latin-ordinary}{c}$
لا يقع في~$\Th{TA}$. فحسبنا إثبات وجود نموذج لـ$\OLId{greek-variable}{Gamma}$.

ولإيجاد نموذج لـ$\OLId{greek-variable}{Gamma}$ نستعمل التراص: نثبت إرضاء كل جزء
منتهٍ من~$\OLId{greek-variable}{Gamma}$، فيلزم إرضاء $\OLId{greek-variable}{Gamma}$ كلها. خذ
$\OLId{greek-variable}{Gamma}_\OLMathNumeral{0} \subseteq \OLId{greek-variable}{Gamma}$ منتهية. وما في $\OLId{greek-variable}{Gamma}_\OLMathNumeral{0}$
إما !!{sentence}s من~$\Th{TA}$، وإما صيغ من الشكل
$\eq/[\OLId{latin-ordinary}{c}][\num{\OLId{latin-ordinary}{n}}]$، وعددها منتهٍ. نجعل $\OLId{latin-ordinary}{k}$ أكبر عدد
يظهر في صيغة $\eq/[\OLId{latin-ordinary}{c}][\num{\OLId{latin-ordinary}{k}}] \in \OLId{greek-variable}{Gamma}_\OLMathNumeral{0}$، ونضع $\OLId{latin-ordinary}{k}=\OLMathNumeral{0}$
إن لم ترد صيغة من هذا الشكل. ثم نبني $\Struct{N_k}$
بتوسيع~$\Struct{N}$ وجعل $\Assign{\OLId{latin-ordinary}{c}}{\OLId{latin-looped}{N}_\OLId{latin-ordinary}{k}}=\OLId{latin-ordinary}{k}+\OLMathNumeral{1}$.
فيكون $\Sat{\OLId{latin-looped}{N}_\OLId{latin-ordinary}{k}}{\OLId{greek-variable}{Gamma}_\OLMathNumeral{0}}$: فإن $!A \in \Th{TA}$
يقتضي $\Sat{\OLId{latin-looped}{N}_\OLId{latin-ordinary}{k}}{!A}$، لأن $\Struct{N_k}$ توافق
$\Struct{N}$ فيما عدا~$\OLId{latin-ordinary}{c}$، ولا يقع $\OLId{latin-ordinary}{c}$ في~$!A$.
وأما $\Sat{\OLId{latin-looped}{N}_\OLId{latin-ordinary}{k}}{\eq/[\OLId{latin-ordinary}{c}][\num{\OLId{latin-ordinary}{n}}]}$ فلأن $\OLId{latin-ordinary}{n} \le \OLId{latin-ordinary}{k}$
و$\Value{\OLId{latin-ordinary}{c}}{\OLId{latin-looped}{N}_\OLId{latin-ordinary}{k}}=\OLId{latin-ordinary}{k}+\OLMathNumeral{1}$. فقد ثبت إرضاء كل جزء منتهٍ،
ومن ثم إرضاء $\OLId{greek-variable}{Gamma}$ بالتراص. واللغة الموسعة قابلة للعد،
فتعطينا مبرهنة لوفنهايم--سكولم النزولية نموذجًا !!{enumerable}؛
ومختزله إلى اللغة الأصلية هو المطلوب.
\end{proof}

\end{document}
